\section{Threats to Validity}
\label{sec:threats}

\subsection{Internal Validity}
\label{sec:threats-internal}

Every \texttt{tabpfn\_*} result in this paper comes from one checkpoint,
verified by SHA-256 hash (Section~\ref{sec:methods-implementation}),
and uses Apple's Metal GPU backend for inference. The whole study,
XGBoost included, ran on one machine
with every library forced to a single thread to avoid a documented
XGBoost/PyTorch conflict. A different checkpoint, backend, or thread
count could shift specific numbers, but there is no reason to expect
the broader pattern, that robustness varies by method and direction,
to depend on any of them. Each (dataset, ratio, contamination rate,
direction) combination was repeated across ten seeds to reduce
dependence on a single train/test split or
corruption realization. Ten repetitions still provide only limited
precision for quantities determined by whether a
seed-averaged value crosses zero, including Correction Harm Rate and
the break-even rate. Additional seeds would further improve
uncertainty characterization of both measures, a limitation
Section~\ref{sec:threats-construct}'s break-even non-monotonicity
finding may partly reflect.

The claim that robustness varies by direction carries its own caveat.
Direction A's and Direction B's contamination rate $c$ are defined against
different reference points (Section~\ref{sec:methods-corruption}), so a
matched nominal $c$ does not correspond to a matched absolute flip count
between the two. Direction A flips a larger fraction of $n_{\min}$'s worth
of rows than Direction B at every tested $c$ above zero. The finding that
direction changes how often a correction fails should be read as a
comparison between two distinct, realistically motivated corruption
processes at their own natural severity levels, not as a
magnitude-controlled manipulation of direction alone. A shared flip-count
definition was considered and rejected, since the two most natural
candidates either require flipping more minority rows than exist at
$r{=}0.05$, or introduce an $r$-dependent mismatch between directions
that grows increasingly severe as $r$ decreases. That mismatch would be
worst precisely at the most imbalanced settings this study emphasizes,
worse than the fixed, $r$-independent gap the current design carries at
every tested $r$.

A related caveat applies to the imbalance ratio $r$ itself. The test
partition is never resampled and keeps each dataset's natural class
distribution (Section~\ref{sec:methods-datasets}), while training is
deliberately pushed to $r$, so the low-$r$ conditions this study
emphasizes combine imbalance with train and test prior shift rather than
isolating imbalance alone. This is the standard evaluation paradigm
across the class-imbalance literature. Training on a rebalanced set and
testing on the natural population is what makes a correction's benefit
deployment-relevant in the first place, but it means this study cannot
cleanly attribute a method's robustness to tolerance of imbalance
specifically as opposed to tolerance of prior shift. This matters most
for \texttt{tabpfn\_threshold}, \texttt{tabpfn\_distpfn}, and
\texttt{tabpfn\_distpfn\_t}, whose underlying mechanism directly targets
prior mismatch (Section~\ref{sec:discussion-simplicity}). A design that
resampled the test partition to match each tested $r$ would remove this
coupling, at the cost of evaluating performance on a population that
does not correspond to any dataset's actual distribution. Consequently,
the present results should be read as robustness under the joint
imbalance/prior-shift setting that rebalanced training and
natural-distribution testing together induce, not as a causal estimate
of imbalance tolerance in isolation.
Section~\ref{sec:results-r050} partially addresses a related but distinct
concern, whether the balanced ratio $r{=}0.50$ (not imbalanced even
before contamination) was itself driving the pooled headline numbers, by
recomputing every contribution metric and statistical test with
$r{=}0.50$ isolated from $r \in \{0.20, 0.10, 0.05\}$. The results
confirm that a substantial share of the DistPFN pair's apparent harm
under Direction B is concentrated at that one ratio. This does not
resolve the train/test prior-shift coupling described above, which
remains present within the $r{\neq}0.50$ subset itself, but it does rule
out $r{=}0.50$ specifically as an unexamined confound in the headline
comparisons.

One correction method's construction interacts with a subset of datasets
in a way worth naming directly. \texttt{xgb\_smote} synthesizes new
minority rows by interpolating in the feature space actually presented to
it, the already-imputed, one-hot-encoded representation
(Section~\ref{sec:methods-datasets}), not the raw categorical values. For
the seven datasets with genuinely categorical columns, this means a
synthetic row can carry a fractional value for what should be a one-hot
indicator, e.g.\ 0.37 rather than 0 or 1, mathematically usable by
XGBoost's tree splits but not a valid category membership.
\texttt{xgb\_ros} is unaffected, since it duplicates real rows rather than
interpolating between them. This is standard practice for applying
off-the-shelf SMOTE to mixed-type data, not an implementation error
specific to this study, and auditing correction methods as they are
commonly deployed, rather than in a bespoke mixed-type-aware variant, is
consistent with this study's stated aim
(Section~\ref{sec:methods-corrections}). Restricting \texttt{xgb\_smote}'s
headline numbers to the 19 datasets with no categorical columns changes
them only slightly: CHR moves from 13.5\%/7.2\% (pooled,
Table~\ref{tab:robustness}) to 13.5\%/6.2\%, and mean CG from
0.023/0.033 to 0.026/0.037, so the categorical datasets are not driving
this method's measured robustness.

\subsection{External Validity}
\label{sec:threats-external}

The 26 datasets audited range from 270 to 60{,}000 raw rows
(Table~\ref{tab:datasets}), but the raw row count is not what any method
actually fits on. The controlled training set every stress cell draws
from is capped by $N_{\text{common}}$ (Section~\ref{sec:methods-ncommon}),
which tops out at 18{,}698 rows even for Adult, the dataset with the
largest minority class in raw terms, because minority availability at
$r{=}0.50$ is the binding constraint regardless of how many majority rows
the raw dataset has. Every condition
tested therefore stays well within TabPFN's intended small-tabular-data
regime but far from the scale where gradient-boosted methods are
typically deployed. Whether the same robustness ranking holds on larger,
higher-dimensional, or more categorical-heavy tabular data, or simply on
a dataset whose minority class alone is large enough to push
$N_{\text{common}}$ higher, is untested. Every dataset audited is also a
binary-classification problem. Whether the same corrections and the same
robustness patterns extend to multi-class imbalance, where a correction
may need to reweight several minority classes at once rather than one, is
untested as well.

The two corruption mechanisms audited are both
class-conditional (Section~\ref{sec:methods-corruption}). The flip
probability depends only on the true class, not on the specific
feature values of the row being flipped. Real-world label noise is
often instance-dependent instead, arising from cases that are
genuinely ambiguous or that share features with the other class, and
this study says nothing about whether the same methods stay robust
under that kind of noise. The four tested contamination rates also stop
at 0.30, with results never extrapolated beyond that point
(Section~\ref{sec:methods-contribution}), so this study cannot say
whether the same patterns hold at higher contamination. Each stress cell
also applies exactly one corruption direction at a time, never a mixture
of both within the same corrupted training set. Real-world label noise
plausibly combines both directions at once, and this study does not test
whether a correction's robustness under one direction in isolation
predicts its robustness once both are present together.

\subsection{Construct Validity}
\label{sec:threats-construct}

Correction Gain, Correction Harm Rate, Regret, and the break-even rate
are original to this paper. They are reported here only on balanced
accuracy, not on the other nine metrics tracked in
Section~\ref{sec:methods-metrics}. A correction ranked robust by CHR
on balanced accuracy is not guaranteed to be equally robust by the
same measure on AUPRC or minority recall, a point the calibration
results in Section~\ref{sec:results-reliability} already illustrate
for a different metric pair. The break-even rate itself assumes a
correction's gain declines monotonically as contamination rises, an
assumption this study's own data does not fully support. Of the 522
(method, dataset, ratio, direction) combinations with a defined
break-even point, 54.6\% show the seed-averaged gain recover to
positive at a higher contamination level after already crossing zero
at a lower one, a majority, not a minority, of defined break-even
points. $c^*$ marks the first noise level a correction stops
paying off, not a guarantee that it stays that way at every higher
level tested. The Prior Sensitivity Gap's oracle
variant assumes access to the true, pre-corruption class prior,
something no real deployment has. It is a diagnostic construct, not a
claim about achievable performance, and Section~\ref{sec:results-rq4}
found it is not even reliably better than the observed-prior variant
it is meant to bound.
